\section*{Guide to the proof}
\addcontentsline{toc}{section}{Guide to the proof}
\label{sec:proof-architecture-v3}

The argument has three stages---location, existence, and amplification---but
the second stage contains most of the new combinatorics. This guide records the
role of each estimate before the detailed proof begins.

\subsection*{1. Locating two first-moment roots}

Let $r_+(n)$ be the continuous root of the ordinary profile first moment, and
let $r_4^{\mathrm{co}}(n)$ be the root of the signed objective restricted to
the four deficits $2,3,4,5$. If $\delta_n$ denotes the phase, put
\[
  A_4(\delta)=\log 2-D_4(\delta),
\]
where $D_4$ is the entropy loss caused by the four-point restriction. The
phase-uniform analysis gives
\[
  r_+(n)-r_4^{\mathrm{co}}(n)
  =
  \left[
    \frac{(\log 2)^2}{4}A_4(\delta_n)+o(1)
  \right]
  \frac{n}{(\log n)^3}.
\]
The signed witness is placed at a tangent-rounded midpoint of these two roots.
The exact conservation correction costs only $O(1)$ classes, so the retained
separation is one half of the leading root gap. Section~11 proves the uniform
certificate
\[
  A_4(\delta)>
  \log\!\left(\frac{1000}{639}\right)
  \qquad(0\le\delta\le1).
\]

\subsection*{2. Extracting the exact overlap structure}

For two ordered signed witnesses, let $r=(r_{ab})$ be their overlap table. The
sum over compatible sign assignments is evaluated before any estimate is
made:
\[
  \operatorname{wt}_{\mathrm{sign}}(r)
  =
  \left(\prod_{a,b}g(r_{ab})\right)2^{\beta(H(r))}.
\]
The product is local; the factor $2^{\beta(H(r))}$ is the cardinality of a
binary cycle space. This identity is the point at which the signed problem
separates into cellwise rewards and a topological residual factor.

The normalized second moment is then split by overlap geometry. Section~7
handles partial diagonals. In every remaining table, a cell above half the
class-size cap is unique in its row and column, so all high cells form a block
matching.

\subsection*{3. Summing high cells before estimating them}

Fix a high-cell support $P$ and realized multiplicities
$j_e=m_e-h_e$. Because $P$ is a matching, the physical choices in different
cells are independent at the level of finite counting. Their complete
aggregate is
\[
  w(P,j)
  =
  \frac{\prod_{e\in P}(s_e)_{j_e}(t_e)_{j_e}}
       {(n)_{\sum_ej_e}\prod_{e\in P}j_e!}
  \prod_{e\in P}g(j_e).
\]
Only after this exact fiber has been summed do we compare $j_e$ with the
full-containment multiplicity $m_e$. The local ratios factor, while the sole
nonlocal change is
\[
  \frac{(n)_{\sum_em_e}}{(n)_{\sum_ej_e}}
  =(n-\textstyle\sum_ej_e)_{\sum_eh_e}
  \le n^{\sum_eh_e}.
\]
Thus the finite-population loss is paid once rather than once per cell.

The positive deficits are summed by a finite optional-choice product. The
full-containment references are then partitioned by endpoint table, and the
resulting table sum is transported to the one-sided partial-diagonal weights
from Section~7. The four-size phase estimates make the total high-cell cost
\[
  \exp\!\left\{o\!\left(\frac{n}{(\log n)^4}\right)\right\}.
\]
The underlying fiber identity, deficit product, and reference regrouping are
valid for any finite endpoint alphabet; only the final transport estimate is
specialized to four consecutive sizes.

\subsection*{4. Restricting residual cycles and amplifying the seed}

After the high cells are exposed, Section~9 assigns one activity $q_{ab}$ to
each residual cell. Expanding the local rewards produces a weighted sum over
even residual edge sets. Deleting the exposed matching is injective on this
family: the symmetric difference of two sets with the same restriction would
be an even subset of a matching, and hence empty. The residual cycle-space sum
is therefore bounded by a full product, giving
\[
  \mathcal A(M,j)
  \le
  \exp\!\left(2\sum_{a,b}q_{ab}\right).
\]
A two-regime estimate makes the exponent
$o(n/(\log n)^4)$ uniformly in the exposed skeleton. Together with the
high-cell estimate, this yields
\[
  \frac{\mathbb E Z^2}{(\mathbb E Z)^2}
  \le
  \exp\!\left\{o\!\left(\frac{n}{(\log n)^4}\right)\right\}.
\]
Paley--Zygmund supplies a witness with probability $e^{-o(n/(\log n)^4)}$.
Section~10 then uses vertex-block bounded differences and a simultaneous
leftover-coloring estimate to turn this rare seed into a high-probability
cocoloring at a cost $o(n/(\log n)^3)$. Intersecting it with the ordinary
chromatic lower bound completes the gap estimate.
